% ═══════════════════════════════════════════════════════════════
% SPANISCH LEHRERHINWEISE LH-05a: LOS HOBBYS
% ═══════════════════════════════════════════════════════════════
% Fachfarbe: #c41e3a (Karminrot)
% ═══════════════════════════════════════════════════════════════

\documentclass[11pt, a4paper]{article}

% ═══════════════════════════════════════════════════════════════
% SW-MODUS TOGGLE
% ═══════════════════════════════════════════════════════════════
\newif\ifswmodus
\swmodusfalse

% ═══════════════════════════════════════════════════════════════
% PAKETE
% ═══════════════════════════════════════════════════════════════

\usepackage[utf8]{inputenc}
\usepackage[T1]{fontenc}
\usepackage[ngerman]{babel}
\usepackage{helvet}
\renewcommand{\familydefault}{\sfdefault}

\usepackage[margin=2cm]{geometry}
\usepackage{enumitem}
\usepackage{tabularx}
\usepackage{booktabs}
\usepackage{longtable}

\usepackage{xcolor}
\usepackage{amssymb}
\usepackage{tcolorbox}
\tcbuselibrary{skins}

\usepackage{fancyhdr}
\usepackage{lastpage}
\usepackage{needspace}

% ═══════════════════════════════════════════════════════════════
% FARBEN
% ═══════════════════════════════════════════════════════════════

\definecolor{spanisch-color}{HTML}{c41e3a}
\definecolor{grau-color}{HTML}{888888}
\definecolor{hellgrau-color}{HTML}{f5f5f5}
\definecolor{basis-color}{HTML}{2a7a4b}
\definecolor{standard-color}{HTML}{2c5aa0}
\definecolor{erweiterung-color}{HTML}{7b1fa2}
\definecolor{loesung-color}{HTML}{c62828}

\definecolor{sw-dark}{HTML}{000000}
\definecolor{sw-gray}{HTML}{666666}
\definecolor{sw-white}{HTML}{ffffff}

\ifswmodus
    \colorlet{spanisch}{sw-dark}
    \colorlet{grau}{sw-gray}
    \colorlet{hellgrau}{sw-white}
    \colorlet{basis}{sw-dark}
    \colorlet{standard}{sw-dark}
    \colorlet{erweiterung}{sw-dark}
    \colorlet{loesung}{sw-dark}
\else
    \colorlet{spanisch}{spanisch-color}
    \colorlet{grau}{grau-color}
    \colorlet{hellgrau}{hellgrau-color}
    \colorlet{basis}{basis-color}
    \colorlet{standard}{standard-color}
    \colorlet{erweiterung}{erweiterung-color}
    \colorlet{loesung}{loesung-color}
\fi

\newcommand{\fachfarbe}{spanisch}

% ═══════════════════════════════════════════════════════════════
% VARIABLEN
% ═══════════════════════════════════════════════════════════════

\newcommand{\thema}{Los hobbys}
\newcommand{\sequenz}{05}
\newcommand{\block}{a}
\newcommand{\fach}{Spanisch}
\newcommand{\klasse}{7}
\newcommand{\dauer}{90 Min. (Doppelstunde)}
\newcommand{\lerngruppengroesse}{24}
\newcommand{\datum}{\today}

% ═══════════════════════════════════════════════════════════════
% KOPF- UND FUSSZEILE
% ═══════════════════════════════════════════════════════════════

\pagestyle{fancy}
\fancyhf{}
\renewcommand{\headrulewidth}{0.4pt}
\renewcommand{\footrulewidth}{0.4pt}

\lhead{\fach{} -- Klasse \klasse}
\chead{\textbf{LH-\sequenz\block{} (Lehrerhinweise)}}
\rhead{\datum}

\lfoot{}
\cfoot{}
\rfoot{Seite \thepage\ von \pageref{LastPage}}

% ═══════════════════════════════════════════════════════════════
% BEFEHLE
% ═══════════════════════════════════════════════════════════════

\newcommand{\B}{\textcolor{basis}{\textbf{[B]}}}
\newcommand{\Std}{\textcolor{standard}{\textbf{[S]}}}
\newcommand{\E}{\textcolor{erweiterung}{\textbf{[E]}}}
\newcommand{\loesung}[1]{\textcolor{loesung}{\textbf{#1}}}
\newcommand{\es}[1]{\textcolor{\fachfarbe}{\textbf{#1}}}

\newtcolorbox{kurzplan}{
    colback=hellgrau,
    colframe=\fachfarbe,
    colbacktitle=white,
    coltitle=\fachfarbe,
    boxrule=1pt,
    arc=0pt,
    fonttitle=\bfseries,
    attach boxed title to top left={yshift=-2mm, xshift=5mm},
    boxed title style={boxrule=0pt, colback=white},
    title=Kurzplan
}

\newtcolorbox{differenzierung}{
    colback=white,
    colframe=grau,
    colbacktitle=white,
    coltitle=grau,
    boxrule=0.5pt,
    arc=0pt,
    fonttitle=\bfseries,
    attach boxed title to top left={yshift=-2mm, xshift=5mm},
    boxed title style={boxrule=0pt, colback=white},
    title=Differenzierung
}

\newtcolorbox{fehlerbox}{
    colback=white,
    colframe=loesung,
    colbacktitle=white,
    coltitle=loesung,
    boxrule=0.5pt,
    arc=0pt,
    fonttitle=\bfseries,
    attach boxed title to top left={yshift=-2mm, xshift=5mm},
    boxed title style={boxrule=0pt, colback=white},
    title=Häufige Fehler
}

\newtcolorbox{materialbox}{
    colback=white,
    colframe=\fachfarbe,
    colbacktitle=white,
    coltitle=\fachfarbe,
    boxrule=0.5pt,
    arc=0pt,
    fonttitle=\bfseries,
    attach boxed title to top left={yshift=-2mm, xshift=5mm},
    boxed title style={boxrule=0pt, colback=white},
    title=Material-Checkliste
}

% ═══════════════════════════════════════════════════════════════
% DOKUMENT
% ═══════════════════════════════════════════════════════════════

\begin{document}

% ═══════════════════════════════════════════════════════════════
% KOPFBEREICH
% ═══════════════════════════════════════════════════════════════

\begin{center}
    {\Large\bfseries\textcolor{\fachfarbe}{Lehrerhinweise: \thema}}\\[0.3em]
    {\normalsize LH-\sequenz\block{} | \dauer}
\end{center}

\vspace{10pt}

% ═══════════════════════════════════════════════════════════════
% 1. KURZPLAN
% ═══════════════════════════════════════════════════════════════

\section*{1. Stundenverlauf (Kurzplan)}

\textbf{1. Stunde (45 Min.)}

\begin{tabularx}{\textwidth}{|c|l|X|l|}
\hline
\textbf{Zeit} & \textbf{Phase} & \textbf{Inhalt} & \textbf{Material} \\
\hline
0--5 & Einstieg & Bildimpuls: Jugendliche bei Freizeitaktivitäten & Tafel/Beamer \\
\hline
5--15 & Input & Vokabeln einführen (9 Hobbys), Aussprache & Bildkarten \\
\hline
15--25 & Übung 1 & Chorisches Sprechen, Einzelkorrektur & -- \\
\hline
25--35 & Übung 2 & Bild-Wort-Zuordnung an Tafel & Wortkarten \\
\hline
35--45 & Sicherung & Tafelanschrieb, ins Heft übertragen & Tafel \\
\hline
\end{tabularx}

\vspace{1em}

\textbf{2. Stunde (45 Min.)}

\begin{tabularx}{\textwidth}{|c|l|X|l|}
\hline
\textbf{Zeit} & \textbf{Phase} & \textbf{Inhalt} & \textbf{Material} \\
\hline
0--5 & Wiederholung & Vokabelquiz: Bild zeigen $\rightarrow$ SuS nennen & Bildkarten \\
\hline
5--15 & Input & ,,¿Qué te gusta hacer?'' + ,,Me gusta...'' & Tafel \\
\hline
15--25 & Anwendung & Partnerarbeit: Dialoge über Hobbys & Dialogkarten \\
\hline
25--35 & Festigung & AB-05a bearbeiten (EA, differenziert) & AB-05a \\
\hline
35--45 & Abschluss & Lösungsbesprechung, HA erteilen & Tafel/ML \\
\hline
\end{tabularx}

% ═══════════════════════════════════════════════════════════════
% 2. DIFFERENZIERUNG
% ═══════════════════════════════════════════════════════════════

\section*{2. Differenzierung}

\begin{differenzierung}
\textbf{Legende:} \B{} Basis (BR) | \Std{} Standard (MR) | \E{} Erweiterung (GY)

\vspace{0.5em}

\textbf{WICHTIG:} Diese Marker sind NUR für die Lehrkraft. Auf Schülermaterial erscheinen KEINE Niveaumarkierungen!
\end{differenzierung}

\vspace{1em}

\begin{tabularx}{\textwidth}{|l|X|X|X|}
\hline
& \B{} \textbf{Basis} & \Std{} \textbf{Standard} & \E{} \textbf{Erweiterung} \\
\hline
\textbf{Aufgabe 1} & Zuordnung mit Bildkarten als Hilfe & Zuordnung selbstständig & + Sätze mit den Vokabeln bilden \\
\hline
\textbf{Aufgabe 2} & Lückentext mit Wortliste & Lückentext selbstständig & + Eigene Kontexte schreiben \\
\hline
\textbf{Aufgabe 3} & Dialog mit Satzanfängen & Dialog mit Stichworten & Freier Dialog + Erweiterung \\
\hline
\textbf{Aufgabe 4} & 1-2 Sätze, kurz & 3 Sätze vollständig & + weitere Hobbys recherchieren \\
\hline
\end{tabularx}

\vspace{1em}

\textbf{Konkrete Hinweise:}
\begin{itemize}
    \item \B{} SuS mit LRS: Zeitzugabe (+10 Min.), vergrößertes AB-Format (A3), Bildkarten zur Unterstützung
    \item \B{} DaZ-SuS: Deutsch-Spanisch-Wortliste vorab besprechen, Tandempartner zuweisen
    \item \E{} Leistungsstarke: Expertenrolle bei Partnerarbeit, 3 weitere Hobbys recherchieren
\end{itemize}

% ═══════════════════════════════════════════════════════════════
% 3. HÄUFIGE FEHLER
% ═══════════════════════════════════════════════════════════════

\section*{3. Häufige Fehler}

\begin{fehlerbox}
\begin{tabularx}{\textwidth}{|l|X|X|}
\hline
\textbf{Fehler} & \textbf{Beispiel} & \textbf{Reaktion} \\
\hline
Aussprache \textit{j} & *[dzh]ugar statt \es{[x]ugar} & Vergleich: dt. ,,ach'', chorisch sprechen \\
\hline
\textit{jugar} ohne \textit{al} & *jugar fútbol & Merkhilfe: Sport braucht ,,al'' \\
\hline
\textit{tocar} mit \textit{al} & *tocar al guitarra & Instrument = \textit{la}, nicht \textit{al} \\
\hline
Aussprache \textit{gu} & *[gu]itarra statt \es{[g]itarra} & gu vor e/i = [g], nicht [gu] \\
\hline
,,Ich mag'' statt ,,Mir gefällt'' & *Yo gusto statt \es{Me gusta} & Struktur: \textit{Me gusta} + Infinitiv \\
\hline
\end{tabularx}
\end{fehlerbox}

% ═══════════════════════════════════════════════════════════════
% 4. MATERIAL-CHECKLISTE
% ═══════════════════════════════════════════════════════════════

\section*{4. Material-Checkliste}

\begin{materialbox}
\begin{itemize}[label=$\square$]
    \item AB-\sequenz\block.pdf (\lerngruppengroesse{} Kopien)
    \item ML-\sequenz\block.pdf (1$\times$ für Lehrkraft)
    \item Lehrwerk: Qué pasa 1, S. 70--72
    \item Bildkarten Hobbys (9 Stück, laminiert, A5)
    \item Wortkarten Hobbys (9 Stück, für Zuordnung an Tafel)
    \item Dialogkarten für Partnerarbeit (12 Paar)
    \item Tafelmarker / Kreide
    \item Optional: LP-\sequenz\block.html (Tablets/PCs für digitalen Lernpfad)
    \item Optional: Audio-Track Unidad 4 (Hörverstehen)
\end{itemize}
\end{materialbox}

% ═══════════════════════════════════════════════════════════════
% 5. LÖSUNGEN
% ═══════════════════════════════════════════════════════════════

\section*{5. Lösungen AB-\sequenz\block}

\textbf{Merkkasten (Übertragung):}

\begin{tabular}{ll}
\es{jugar al fútbol} & $\rightarrow$ \loesung{Fußball spielen} \\
\es{escuchar música} & $\rightarrow$ \loesung{Musik hören} \\
\es{ver la tele} & $\rightarrow$ \loesung{fernsehen} \\
\es{leer} & $\rightarrow$ \loesung{lesen} \\
\es{nadar} & $\rightarrow$ \loesung{schwimmen} \\
\es{montar en bici} & $\rightarrow$ \loesung{Fahrrad fahren} \\
\es{tocar la guitarra} & $\rightarrow$ \loesung{Gitarre spielen} \\
\es{chatear} & $\rightarrow$ \loesung{chatten} \\
\es{jugar a videojuegos} & $\rightarrow$ \loesung{Videospiele spielen} \\
\end{tabular}

\vspace{1em}

\textbf{Aufgabe 1: Zuordnung} (4 Punkte)

\begin{tabular}{ll}
\es{jugar al fútbol} & $\rightarrow$ \loesung{Fußball spielen} \\
\es{escuchar música} & $\rightarrow$ \loesung{Musik hören} \\
\es{ver la tele} & $\rightarrow$ \loesung{fernsehen} \\
\es{tocar la guitarra} & $\rightarrow$ \loesung{Gitarre spielen} \\
\end{tabular}

\vspace{1em}

\textbf{Aufgabe 2: Lückentext} (5 Punkte)

\begin{tabular}{ll}
a) & \loesung{nadar} \\
b) & \loesung{montar en bici} \\
c) & \loesung{chatear} \\
d) & \loesung{jugar a videojuegos} \\
e) & \loesung{leer} \\
\end{tabular}

\vspace{1em}

\textbf{Aufgabe 3: Dialog} (5 Punkte)

Mögliche Antworten:
\begin{itemize}
    \item B1: \loesung{Me llamo [Name]. / Soy [Name].}
    \item B2: \loesung{Me gusta [Hobby].} (z.B. Me gusta nadar.)
    \item B3: \loesung{Sí, me gusta nadar.} oder \loesung{No, no me gusta nadar.}
    \item B4: \loesung{Me gusta también [Hobby].}
    \item B5: \loesung{¡Hasta luego! / ¡Adiós!}
\end{itemize}

\textit{Bewertung: Je 1P pro korrekte, sinnvolle Antwort}

\vspace{1em}

\textbf{Aufgabe 4: Freie Produktion} (6 Punkte)

Bewertungskriterien pro Satz (max. 2P):
\begin{itemize}
    \item Struktur ,,Me gusta + Infinitiv'' korrekt: 1P
    \item Hobby korrekt geschrieben: 1P
\end{itemize}

\textit{Hinweis: Verschiedene Hobbys verwenden! 3 Sätze = max. 6P}

% ═══════════════════════════════════════════════════════════════
% 6. HAUSAUFGABE
% ═══════════════════════════════════════════════════════════════

\section*{6. Hausaufgabe}

\begin{itemize}
    \item Lehrwerk S. 71, Aufgabe 3 (Vokabelübung)
    \item Vokabeln lernen: \textit{jugar al fútbol} bis \textit{jugar a videojuegos} (9 Wörter)
    \item Optional: Lernpfad LP-\sequenz\block.html (digital)
    \item \E{} Zusatz: 3 weitere Hobbys recherchieren und aufschreiben
\end{itemize}

% ═══════════════════════════════════════════════════════════════
% 7. REFLEXIONSNOTIZEN
% ═══════════════════════════════════════════════════════════════

\section*{7. Reflexionsnotizen (nach Durchführung)}

\begin{tabularx}{\textwidth}{|l|X|}
\hline
\textbf{Was lief gut?} & \\[2em]
\hline
\textbf{Was lief nicht?} & \\[2em]
\hline
\textbf{Zeitmanagement} & \\[2em]
\hline
\textbf{Für nächstes Mal} & \\[2em]
\hline
\end{tabularx}

\end{document}
